% ---------------------------------------------------------------------
% LaTeX (Overleaf) Курсовой проект
% Тема: Разработка и реализация мультимодальной агентной системы Chronobiotic
% ---------------------------------------------------------------------
\UseRawInputEncoding
\documentclass[a4paper,14pt]{extarticle}
\usepackage[utf8]{inputenc}
\usepackage[T2A]{fontenc}
\usepackage[russian]{babel}
\usepackage{csquotes}

% Поля: 3см слева, 1.5см справа, 2см сверху/снизу
\usepackage{geometry}
\geometry{
	a4paper,
	left=3cm,
	right=1.5cm,
	top=2cm,
	bottom=2cm
}

% Настройка межстрочного интервала (1.15)
\usepackage{setspace}
\setstretch{1.15}

% Настройка отступа абзаца (1 см) для всего документа
\usepackage{indentfirst}
\setlength{\parindent}{1cm}

% Выравнивание по ширине для всего документа
\usepackage{ragged2e}
\justifying

% ПОЛНОСТЬЮ ПЕРЕОПРЕДЕЛЯЕМ ЗАГОЛОВКИ
\usepackage{titlesec}

% Глава 1 уровня (section) - ЖИРНЫЙ, с отступом
\titleformat{\section}[block]{\normalfont\bfseries\fontsize{16pt}{19pt}\selectfont}{\hspace{1cm}\thesection}{0.5em}{}

% Глава 2 уровня (subsection) - ЖИРНЫЙ, с отступом
\titleformat{\subsection}[block]{\normalfont\bfseries\fontsize{14pt}{17pt}\selectfont}{\hspace{1cm}\thesubsection}{0.5em}{}

% Глава 3 уровня (subsubsection) - ЖИРНЫЙ, с отступом
\titleformat{\subsubsection}[block]{\normalfont\bfseries\fontsize{14pt}{17pt}\selectfont}{\hspace{1cm}\thesubsubsection}{0.5em}{}

% Для ненумерованных заголовков (Введение, Заключение, Список литературы, Приложения)
\titleformat{name=\section,numberless}[block]{\normalfont\bfseries\fontsize{16pt}{19pt}\selectfont}{\hspace{1cm}}{0em}{}

% Отступы вокруг заголовков
\titlespacing*{\section}{0pt}{18pt}{12pt}
\titlespacing*{\subsection}{0pt}{12pt}{6pt}
\titlespacing*{\subsubsection}{0pt}{12pt}{6pt}
\titlespacing*{name=\section,numberless}{0pt}{18pt}{12pt}

% Пакет для колонтитулов
\usepackage{fancyhdr}
\pagestyle{fancy}
\fancyhf{}
\fancyfoot[C]{\thepage}
\renewcommand{\headrulewidth}{0pt}
\setlength{\footskip}{1.5cm}

% Пакет для листингов кода
\usepackage{listings}
\usepackage{xcolor}
\definecolor{codegray}{gray}{0.95}
\lstdefinestyle{mystyle}{
	backgroundcolor=\color{codegray},
	basicstyle=\ttfamily\fontsize{13pt}{15pt}\selectfont,
	breakatwhitespace=false,
	breaklines=true,
	captionpos=t,
	keepspaces=true,
	numbers=none,
	showspaces=false,
	showstringspaces=false,
	showtabs=false,
	tabsize=2,
	frame=single,
	frameround=tttt,
	framesep=5pt,
	aboveskip=12pt,
	belowskip=12pt
}
\lstset{style=mystyle}
\renewcommand{\lstlistingname}{Листинг}

% Пакет для подписей рисунков и таблиц
\usepackage{caption}
\captionsetup[figure]{
	name={Рисунок},
	labelsep=space,
	justification=centering,
	font={normalsize},
	skip=6pt,
	position=below
}
\captionsetup[table]{
	name={Таблица},
	labelsep=space,
	justification=centering,
	font={normalsize},
	skip=6pt,
	position=above
}
\captionsetup[lstlisting]{
	labelsep=space,
	justification=centering,
	font={normalsize},
	skip=6pt,
	position=above
}

% Пакет для таблиц
\usepackage{booktabs}
\usepackage{array}

% Пакет для графики
\usepackage{graphicx}
\graphicspath{{images/}}

% Пакет для ссылок
\usepackage[colorlinks=false, hidelinks]{hyperref}

% Настройка списков
\usepackage{enumitem}
\setlist{
	leftmargin=2.0cm,
	itemindent=0.2cm,
	labelwidth=0.8cm,
	labelsep=0.3cm,
	align=left,
	topsep=6pt,
	partopsep=0pt,
	parsep=0pt,
	itemsep=3pt
}

\setlist[itemize,1]{leftmargin=2.2cm, label=--}
\setlist[itemize,2]{leftmargin=2.7cm, label=$\circ$}
\setlist[enumerate,1]{leftmargin=2.2cm}
\setlist[enumerate,2]{leftmargin=2.7cm}

% Отключение автоматической расстановки переносов
\tolerance=1
\emergencystretch=\maxdimen
\hyphenpenalty=10000
\hbadness=10000

% Отключение разреженного/уплотненного текста
\spaceskip=1.0\fontdimen2\font plus 1.0\fontdimen3\font minus 1.0\fontdimen4\font

% Начало документа
\begin{document}
	
	% ---------------------------------------------------------------------
	% ТИТУЛЬНЫЙ ЛИСТ
	% ---------------------------------------------------------------------
	\begin{titlepage}
		\begin{center}
			\large{Минобрнауки России}
			
			\large{Федеральное государственное бюджетное образовательное учреждение высшего образования\\
				«Сыктывкарский государственный университет имени Питирима Сорокина»\\
				(ФГБОУ ВО «СГУ им. Питирима Сорокина»)}
			
			\large{Институт точных наук и информационных технологий}
			
			\large{Кафедра прикладной математики и компьютерных наук}
			
			\vspace{4cm}
			
			\textbf{\Large{КУРСОВОЙ ПРОЕКТ}}
			
			\large{по дисциплине «Проектирование в профессиональной сфере»}
			
			\vspace{1cm}
			
			\textbf{\Large{Разработка и реализация мультимодальной агентной системы Chronobiotic для интеллектуального анализа данных в области хронобиологии}}
			
			\vspace{1.5cm}
			
			\begin{flushleft}
				\large{Направление подготовки:\\
					01.03.02 Прикладная математика и информатика\\
					Направленность (профиль программы):\\
					Искусственный интеллект}
			\end{flushleft}
			
			\vspace{2cm}
			
			\begin{flushleft}
				\large{Исполнитель: Буренкова Полина Алексеевна 131-ПМо\\
					\underline{\hspace{5cm}}}
				
				\large{Научный руководитель: доцент, к.ф.-м.н. Котелина Надежда Олеговна,\\
					\underline{\hspace{5cm}}}
				
				\vspace{1.5cm}
				
				\large{Руководитель ОПОП:\\
					зав. кафедрой ПМиКН, к.ф.-м.н., доцент\\
					Ермоленко А. В. \underline{\hspace{3cm}}}
			\end{flushleft}
			
			\vfill
			
			\large{Сыктывкар\\2026}
		\end{center}
	\end{titlepage}
	
	% ---------------------------------------------------------------------
	% СОДЕРЖАНИЕ
	% ---------------------------------------------------------------------
	\setcounter{page}{2}
	\tableofcontents
	\newpage
	
	% ---------------------------------------------------------------------
	% ВВЕДЕНИЕ
	% ---------------------------------------------------------------------
	\section*{Введение}
	\addcontentsline{toc}{section}{Введение}
	
	\textbf{Актуальность исследования.} Современный этап развития биоинформатики и медицинской химии характеризуется экспоненциальным ростом объема научных данных, что требует разработки новых методов их интеллектуальной обработки. Особую значимость приобретают специализированные информационные системы в области хронобиологии – науки о биологических ритмах и их влиянии на физиологические процессы. Хронобиотики, как класс соединений, модулирующих циркадные ритмы, представляют значительный интерес для фундаментальной науки и практической медицины. Однако существующие базы данных хронобиотиков, включая разработанную ранее базу объемом 340 записей, обладают ограниченным функционалом интеллектуального анализа и не поддерживают современные способы взаимодействия с пользователем. Создание агентной системы с мультимодальным интерфейсом, интегрирующей методы Retrieval-Augmented Generation (RAG), Knowledge-Augmented Generation (KAG) и большие языковые модели (LLM), позволит существенно повысить эффективность работы исследователей с данными в области хронобиологии.
	
	\textbf{Степень разработанности проблемы.} В настоящее время существует значительное количество исследований, посвященных применению агентных систем в биоинформатике. Работы таких авторов, как Wooldridge, Jennings, Ferber, заложили теоретические основы многоагентных систем. Исследования Lewis, Guu, Karpukhin в области RAG-архитектур, а также работы Ji, Pan, Wang по графовым базам знаний определили современные подходы к гибридному поиску информации. Развитие LLM, включая открытую модель BLOOM (BigScience Large Open-science Open-access Multilingual), создало предпосылки для создания специализированных медицинских систем. Однако комплексные решения, объединяющие аналитические, голосовые и мультиязычные агенты для работы с узкоспециализированными базами данных хронобиотиков, в научной литературе представлены недостаточно.
	
	\textbf{Цель и задачи исследования.} Целью настоящей работы является разработка и реализация мультимодальной агентной системы Chronobiotic, обеспечивающей интеллектуальный анализ базы данных хронобиотиков с использованием гибридных методов поиска и больших языковых моделей.
	
	Для достижения поставленной цели необходимо решить следующие задачи:
	\begin{enumerate}
		\item Провести анализ современных подходов к построению агентных систем и методов обработки специализированных медицинских данных.
		\item Исследовать архитектурные решения для интеграции LLM (BLOOM) с реляционными базами данных.
		\item Разработать классификацию и архитектуру специализированных агентов для анализа хронобиотиков.
		\item Реализовать систему агентов, включающую аналитические, хронобиологические, голосовые и мультиязычные компоненты.
		\item Создать гибридную систему поиска на основе комбинации RAG и KAG методов поверх PostgreSQL базы данных.
		\item Интегрировать голосовые интерфейсы (STT/TTS) и механизмы мультиязычной поддержки.
		\item Провести fine-tuning модели BLOOM на предметно-ориентированных данных хронобиотиков.
		\item Осуществить тестирование разработанной системы и оценку эффективности ее компонентов.
		\item Разработать пользовательский интерфейс с поддержкой мультимодального взаимодействия.
	\end{enumerate}
	
	\textbf{Объект и предмет исследования.} Объектом исследования являются процессы интеллектуального анализа и извлечения информации из специализированных медицинских баз данных. Предметом исследования выступают методы, алгоритмы и архитектурные решения для построения мультимодальных агентных систем с интеграцией больших языковых моделей.
	
	\textbf{Научная новизна работы} заключается в:
	\begin{itemize}
		\item разработке архитектуры комплексной мультимодальной агентной системы, интегрирующей аналитические, голосовые и мультиязычные компоненты для работы с предметно-ориентированной базой данных;
		\item создании гибридной системы поиска, комбинирующей RAG и KAG методы поверх реляционной базы данных PostgreSQL без необходимости миграции в специализированные графовые хранилища;
		\item адаптации модели BLOOM к предметной области хронобиологии посредством fine-tuning на ограниченном наборе данных (340 записей) с использованием методов few-shot обучения;
		\item разработке методологии координации специализированных агентов для анализа медицинских данных с учетом их мультимодальной природы.
	\end{itemize}
	
	\textbf{Теоретическая значимость работы} определяется вкладом в развитие методов построения специализированных агентных систем для медицинских приложений, а также в исследование возможностей комбинирования различных архитектурных подходов (RAG, KAG, LLM) для работы с ограниченными предметно-ориентированными данными.
	
	\textbf{Практическая значимость работы} заключается в создании функционирующей агентной системы Chronobiotic, которая может быть использована исследователями в области хронобиологии и хронофармакологии для эффективного поиска и анализа информации о биологически активных соединениях, модулирующих циркадные ритмы.
	
	\textbf{Методология и методы исследования.} В работе используются методы системного анализа, объектно-ориентированного проектирования, теории агентных систем, методы обработки естественного языка, машинного обучения, технологии реляционных и графовых баз данных, методы оценки качества информационных систем.
	
	\textbf{Положения, выносимые на защиту:}
	\begin{enumerate}
		\item Архитектура мультимодальной агентной системы, обеспечивающая координацию специализированных агентов различного типа при анализе данных в области хронобиологии.
		\item Метод гибридного поиска информации, основанный на комбинации векторного представления (RAG) и графовых структур (KAG) поверх реляционной базы данных.
		\item Результаты fine-tuning модели BLOOM на ограниченном наборе предметно-ориентированных данных (340 записей).
		\item Программная реализация агентной системы Chronobiotic с поддержкой голосовых и мультиязычных интерфейсов.
	\end{enumerate}
	
	\textbf{Достоверность и обоснованность результатов} обеспечивается корректным применением современных методов проектирования и разработки программных систем, экспериментальной проверкой разработанных алгоритмов на реальных данных, использованием общепризнанных метрик оценки качества информационных систем.
	
	\textbf{Апробация работы.} Основные положения и результаты работы будут представлены на научных семинарах кафедры, а также планируются к публикации в сборниках научных трудов и выступлениям на конференциях соответствующего профиля.
	
	\textbf{Структура и объем работы.} Работа состоит из введения, трех глав, заключения, списка литературы (70 наименований) и приложений. Общий объем работы составляет 60-70 страниц машинописного текста.
	\newpage
	
	% ---------------------------------------------------------------------
	% ГЛАВА 1
	% ---------------------------------------------------------------------
	\section{Теоретические основы построения мультимодальных агентных систем для анализа медико-биологических данных}
	
	\subsection{Эволюция и современное состояние исследований в области интеллектуальных агентных систем}
	
	\subsubsection{Генезис концепции интеллектуальных агентов: от экспертных систем к мультиагентным архитектурам}
	История развития интеллектуальных систем берет начало от экспертных систем 1970-х годов, которые оперировали статическими базами знаний и правилами логического вывода. Однако ограниченность таких систем в решении динамических и распределенных задач привела к появлению концепции интеллектуальных агентов в 1990-х годах. Работы М. Вулдриджа и Н. Дженнингса заложили фундаментальные основы теории агентов, определив такие ключевые свойства, как автономность, реактивность, проактивность и социальное поведение. Эволюция от отдельных агентов к мультиагентным системам позволила решать сложные задачи путем декомпозиции на подзадачи, распределенные между специализированными агентами.
	
	\subsubsection{Типология и классификация агентов в современных информационных системах}
	В современной теории агентных систем существует множество классификаций агентов по различным признакам. По типу архитектуры выделяют рефлексивные (действующие по правилам), когнитивные (обладающие моделью мира) и гибридные агенты. По выполняемым функциям различают информационные, аналитические, коммуникационные и управляющие агенты. В контексте данной работы особый интерес представляют мультимодальные агенты, способные обрабатывать разнородные типы данных (текст, речь, химические структуры).
	
	\subsubsection{Архитектурные паттерны построения многоагентных систем}
	Анализ литературы позволяет выделить несколько ключевых архитектурных паттернов: горизонтальная архитектура (агенты взаимодействуют напрямую), вертикальная (иерархическая) и архитектура на основе доски объявлений (blackboard). Для разрабатываемой системы Chronobiotic наиболее перспективной представляется гибридная архитектура с центральным менеджером агентов, обеспечивающим координацию, и прямым взаимодействием агентов для решения конкретных задач.
	
	\subsection{Применение агентных систем в биоинформатике и медицинских исследованиях}
	
	\subsubsection{Специфика задач биоинформатики и требования к агентным системам}
	Задачи биоинформатики характеризуются гетерогенностью данных, необходимостью интеграции разнородных источников информации (базы данных, научная литература, экспериментальные данные), высокой размерностью и сложностью взаимосвязей. Агентные системы для биоинформатики должны обеспечивать масштабируемость, интероперабельность, возможность работы с неполными и противоречивыми данными.
	
	\subsubsection{Обзор существующих агентных систем в области медицинской химии и фармакологии}
	Существующие системы, такие как ChemAgen, BioAgent, PharmaS, демонстрируют эффективность агентного подхода для решения задач виртуального скрининга, прогнозирования свойств соединений, поиска лекарственных мишеней. Однако большинство из них ориентированы на узкий класс задач и не предоставляют мультимодальных интерфейсов взаимодействия.
	
	\subsubsection{Проблемы и ограничения современных медицинских агентных систем}
	К основным проблемам относятся сложность интеграции с существующими базами данных, недостаточная поддержка естественно-языковых интерфейсов, ограниченные возможности по работе с новейшими моделями машинного обучения, включая LLM. Отсутствие гибких механизмов адаптации к новым предметным областям также сдерживает их широкое применение.
	
	\subsection{Большие языковые модели: архитектура, методы обучения и области применения}
	
	\subsubsection{Архитектурные особенности трансформерных моделей}
	Трансформерная архитектура, предложенная в работе Vaswani et al. ``Attention Is All You Need'', основана на механизме самовнимания, позволяющем модели учитывать контекстные взаимосвязи между всеми элементами последовательности. Многослойные кодеры и декодеры, предобученные на огромных корпусах текстов, формируют языковые модели, способные решать широкий спектр задач NLP.
	
	\subsubsection{Модель BLOOM: архитектура, мультиязычность, открытый доступ}
	Модель BLOOM (BigScience Large Open-science Open-access Multilingual Language Model) является одной из наиболее значимых открытых моделей. Обученная на 46 естественных языках и 13 программных языках, она демонстрирует высокие результаты в задачах понимания и генерации текста. Открытый доступ и возможность дообучения делают ее идеальной платформой для создания специализированных систем в области хронобиологии.
	
	\subsubsection{Методы fine-tuning языковых моделей для предметных областей}
	Для адаптации LLM к конкретной предметной области используются методы дообучения. Полное дообучение всех параметров требует значительных вычислительных ресурсов. Альтернативой являются методы Parameter-Efficient Fine-Tuning (PEFT), такие как LoRA (Low-Rank Adaptation), которые позволяют дообучать лишь небольшую часть параметров, сохраняя высокое качество адаптации.
	
	\subsubsection{Проблемы адаптации LLM к работе с ограниченными наборами данных}
	Работа с ограниченными предметно-ориентированными данными (в данном случае 340 записей) представляет особую сложность из-за риска переобучения и потери обобщающей способности. Методы few-shot обучения и аугментации данных позволяют частично преодолеть эти ограничения.
	
	\subsection{Гибридные методы извлечения информации: RAG и KAG}
	
	\subsubsection{Архитектура Retrieval-Augmented Generation (RAG): принципы и компоненты}
	RAG-архитектура объединяет этапы поиска информации и генерации ответа. Компонент ретривера находит релевантные фрагменты в базе знаний, которые затем подаются в LLM в качестве контекста для генерации ответа. Это позволяет модели оперировать актуальными и верифицированными данными, снижая риск галлюцинаций.
	
	\subsubsection{Knowledge-Augmented Generation (KAG): графы знаний как источник информации}
	KAG расширяет концепцию RAG, используя структурированные графы знаний. Вместо простых текстовых фрагментов модель получает доступ к узлам и связям графа, что обеспечивает более глубокое понимание взаимосвязей между сущностями.
	
	\subsubsection{Сравнительный анализ подходов и возможности их комбинирования}
	RAG эффективен для поиска по неструктурированным текстам, KAG – для ответов на вопросы, требующие понимания сложных взаимосвязей. Их комбинирование (гибридный подход) позволяет использовать преимущества обоих методов, обеспечивая полноту и точность ответов.
	
	\subsubsection{Особенности реализации гибридного поиска в реляционных базах данных}
	Реализация гибридного поиска поверх реляционной СУБД, такой как PostgreSQL, требует расширения ее функциональности для поддержки векторного поиска (расширение pgvector) и построения графовых структур на основе реляционных таблиц.
	
	\subsection{Технологии обработки речи и мультиязычные системы в научных приложениях}
	
	\subsubsection{Современные методы распознавания речи (STT) для научных приложений}
	Модель Whisper от OpenAI демонстрирует высокое качество распознавания русскоязычной речи, включая научную терминологию, что делает ее предпочтительным выбором для системы Chronobiotic.
	
	\subsubsection{Технологии синтеза речи (TTS): сравнительный анализ подходов}
	Современные TTS-системы, такие как Tacotron, WaveNet, VITS, обеспечивают высокое качество синтеза. Выбор конкретного решения определяется требованиями к реалистичности голоса, задержкам и доступности для русского языка.
	
	\subsubsection{Методы автоматического определения языка и локализации контента}
	Автоматическое определение языка входящего запроса и последующая локализация ответа являются ключевыми для обеспечения мультиязычного диалога. Используются как классификаторы на основе n-грамм, так и нейросетевые подходы.
	
	\subsubsection{Особенности обработки научной терминологии в голосовых интерфейсах}
	Научная терминология, включающая латинские названия, химические формулы, требует специальной обработки в системах STT/TTS. Создание специализированных словарей и аугментация обучающих данных позволяют повысить точность распознавания.
	
	\subsection{Хронобиология как предметная область: анализ данных и информационные ресурсы}
	
	\subsubsection{Теоретические основы хронобиологии и хронофармакологии}
	Хронобиология изучает биологические ритмы, их механизмы и значение для жизнедеятельности организмов. Ключевым понятием являются циркадные ритмы – эндогенные колебания физиологических процессов с периодом около 24 часов. Хронофармакология исследует влияние времени введения лекарств на их эффективность и безопасность.
	
	\subsubsection{Обзор существующих баз данных по хронобиотикам и родственным соединениям}
	Публичные базы данных, такие как CircaDB, BioGPS, содержат информацию об экспрессии генов циркадных часов. Однако специализированные базы хронобиотиков – веществ, модулирующих циркадные ритмы, – в открытом доступе практически отсутствуют.
	
	\subsubsection{Проблемы информационного обеспечения исследований в области хронобиологии}
	Основными проблемами являются разрозненность данных, отсутствие стандартизированных форматов представления, сложность интеграции информации о структуре веществ и их биологических эффектах.
	
	\subsubsection{Структура и особенности разработанной базы данных хронобиотиков (340 записей)}
	В рамках предшествующих исследований была разработана база данных, содержащая 340 записей о хронобиотиках. Структура включает информацию о химическом названии, структуре (SMILES), молекулярных дескрипторах, типе биологической активности, целевых генах/белках, литературных источниках. Данный объем данных, с одной стороны, является ограниченным для обучения LLM ``с нуля'', с другой – достаточным для дообучения и создания эффективной системы поиска.
	
	\subsection{Методология оценки эффективности интеллектуальных информационных систем}
	
	\subsubsection{Критерии и метрики качества для агентных систем}
	Для оценки агентных систем используются как общесистемные метрики (производительность, масштабируемость, отказоустойчивость), так и специфические: время реакции агента, точность выполнения задач, полнота охвата информации.
	
	\subsubsection{Методы оценки точности извлечения информации и генерации ответов}
	Для оценки качества извлечения информации применяются метрики Precision, Recall, F1-мера. Для генеративных ответов используются метрики BLEU, ROUGE, METEOR, а также экспертная оценка специалистами предметной области.
	
	\subsubsection{Подходы к оценке пользовательского опыта в мультимодальных системах}
	Оценка пользовательского опыта (UX) включает анкетирование, анализ поведенческих метрик, проведение интервью. Специфика мультимодальных систем требует оценки удобства переключения между модальностями (текст, голос).
	
	\subsubsection{Стандарты и методологии тестирования медицинских информационных систем}
	Разработка систем для медицинских приложений должна учитывать стандарты качества и безопасности, такие как ISO 13485, IEC 62304, рекомендации FDA по программному обеспечению как медицинскому изделию.
	
	\subsection*{Выводы по главе 1}
	\addcontentsline{toc}{subsection}{Выводы по главе 1}
	
	В первой главе проведен комплексный анализ современного состояния исследований в области интеллектуальных агентных систем, больших языковых моделей, гибридных методов поиска информации и технологий мультимодальных интерфейсов. Установлено, что, несмотря на развитие отдельных технологий, отсутствуют комплексные решения, интегрирующие аналитические, голосовые и мультиязычные компоненты для работы с узкоспециализированными базами данных в области хронобиологии. Обоснован выбор архитектуры многоагентной системы и открытой языковой модели BLOOM в качестве основы для разработки. Выявлена необходимость применения гибридных методов поиска (RAG+KAG) для эффективной работы с гетерогенными данными. Сформулированы требования к разрабатываемой системе и определены критерии оценки ее эффективности.
	
	\newpage
	
	% ---------------------------------------------------------------------
	% ГЛАВА 2
	% ---------------------------------------------------------------------
	\section{Проектирование мультимодальной агентной системы Chronobiotic}
	
	\subsection{Системный анализ предметной области и спецификация требований}
	
	\subsubsection{Функциональные требования к системе на основе анализа пользовательских сценариев}
	На основе анализа пользовательских сценариев (исследователь-хронобиолог, студент, фармаколог) сформулированы следующие функциональные требования:
	\begin{itemize}
		\item Система должна принимать запросы пользователя на естественном языке (текстовый и голосовой ввод).
		\item Система должна осуществлять поиск информации в базе данных хронобиотиков по названию вещества, химической структуре (SMILES), молекулярным свойствам.
		\item Система должна извлекать информацию из загруженных пользователем научных статей (PDF) и интегрировать ее с существующей базой знаний.
		\item Система должна предоставлять структурированный ответ с указанием источников информации и возможностью детализации.
		\item Система должна поддерживать мультиязычный диалог (русский, английский).
	\end{itemize}
	
	\subsubsection{Нефункциональные требования и ограничения системы}
	К нефункциональным требованиям относятся:
	\begin{itemize}
		\item Время ответа на типовой запрос не должно превышать 10 секунд.
		\item Система должна быть модульной и масштабируемой.
		\item Обеспечение безопасности данных и разграничения доступа.
		\item Возможность развертывания как на локальном сервере с GPU, так и в облачной инфраструктуре.
		\item Интеграция с СУБД PostgreSQL, содержащей данные о 340 хронобиотиках.
	\end{itemize}
	
	\subsubsection{Анализ структуры и объема существующей базы данных хронобиотиков}
	База данных включает следующие основные таблицы:
	\begin{itemize}
		\item \texttt{substances} – идентификатор, название, химическая формула, SMILES-представление.
		\item \texttt{properties} – молекулярная масса, logP, количество доноров/акцепторов водородной связи.
		\item \texttt{bio\_activity} – тип активности, количественные показатели (IC50, EC50), целевой ген/белок.
		\item \texttt{literature\_sources} – библиографические ссылки на статьи.
		\item \texttt{circadian\_effects} – влияние на фазу, амплитуду, период циркадных ритмов.
	\end{itemize}
	Объем данных – 340 записей о веществах, что является достаточным для создания эффективной системы поиска и дообучения LLM.
	
	\subsubsection{Определение границ системы и внешних интерфейсов}
	Система Chronobiotic взаимодействует со следующими внешними компонентами:
	\begin{itemize}
		\item Пользователь (через веб-интерфейс и голосовой интерфейс).
		\item Внешние базы данных (PubChem, ChEMBL) для обогащения информации (опционально).
		\item Системы машинного перевода (при необходимости).
	\end{itemize}
	
	\subsection{Архитектура комплексной мультимодальной агентной системы}
	
	\subsubsection{Концептуальная архитектура системы и принципы взаимодействия компонентов}
	Архитектура системы является модульной и строится вокруг ядра – Менеджера агентов. Менеджер принимает пользовательский запрос, определяет его тип и делегирует выполнение соответствующим специализированным агентам. Взаимодействие между агентами осуществляется через асинхронную систему сообщений.
	
	\subsubsection{Ядро агентной системы: менеджер агентов и механизмы координации}
	Менеджер агентов реализует следующие функции:
	\begin{itemize}
		\item Регистрация и дерегистрация агентов.
		\item Маршрутизация запросов на основе анализа компетенций агентов.
		\item Координация выполнения сложных запросов, требующих участия нескольких агентов.
		\item Мониторинг состояния агентов и обработка сбоев.
	\end{itemize}
	
	\subsubsection{Организация коммуникации между агентами различных типов}
	Для коммуникации используется асинхронная очередь сообщений (RabbitMQ/Kafka). Каждый агент имеет уникальный идентификатор и подписывается на типы сообщений, которые он способен обрабатывать.
	
	\subsubsection{Интеграция с внешними системами и источниками данных}
	Интеграция с внешними источниками (PubChem) осуществляется через REST API. Для работы с PDF-документами используются библиотеки для извлечения текста (PyPDF2, pdfplumber).
	
	\subsection{Проектирование аналитических агентов для химического и фармакологического анализа}
	
	\subsubsection{Агент химического анализа (Chemical Analyzer Agent)}
	Агент отвечает за проверку валидности SMILES-представлений, расчет молекулярных дескрипторов (с использованием RDKit), идентификацию функциональных групп.
	
	\subsubsection{Агент оценки токсичности (Toxicity Estimator Agent)}
	Агент использует предобученные модели (например, на основе данных Tox21) для прогнозирования потенциальной токсичности веществ.
	
	\subsubsection{Агент анализа эффективности (Efficacy Evaluator Agent)}
	Анализирует данные о биологической активности (IC50, EC50) и оценивает потенциальную эффективность вещества как модулятора циркадных ритмов.
	
	\subsubsection{Агент анализа взаимодействий (Interaction Analyzer Agent)}
	Агент исследует взаимодействия веществ с белками-мишенями, используя данные из базы и литературных источников.
	
	\subsubsection{Агент предсказания свойств (Property Predictor Agent)}
	Прогнозирует физико-химические и биологические свойства веществ на основе их структуры с использованием методов машинного обучения.
	
	\subsubsection{Агент поиска аналогий (Similarity Finder Agent)}
	Реализует поиск структурных аналогов заданного вещества по базе данных с использованием различных мер подобия (Tanimoto, Dice).
	
	\subsection{Проектирование предметно-ориентированных хронобиологических агентов}
	
	\subsubsection{Экспертный агент по хронобиологии (Chronobiology Expert Agent)}
	Содержит базу знаний о циркадных ритмах, генах часов (CLOCK, BMAL1, PER, CRY), механизмах их регуляции. Используется для интерпретации результатов анализа.
	
	\subsubsection{Агент поиска в базе хронобиотиков (Chronobiotics Searcher Agent)}
	Основной агент для выполнения поисковых запросов по базе данных хронобиотиков.
	
	\subsubsection{Агент анализа клинических данных (Clinical Data Finder Agent)}
	Осуществляет поиск информации о клинических исследованиях хронобиотиков в открытых источниках (ClinicalTrials.gov).
	
	\subsubsection{Агент извлечения информации из литературы (Literature Miner Agent)}
	Извлекает информацию из научных статей (PDF), используя методы NLP и fine-tuned NER-модели для идентификации сущностей (вещества, гены, эффекты).
	
	\subsubsection{Агент исследования механизмов действия (Mechanism Researcher Agent)}
	Анализирует данные о сигнальных путях и механизмах действия веществ.
	
	\subsubsection{Агент комплексного анализа веществ (Substance Analyzer Agent)}
	Координирует работу других аналитических агентов для проведения комплексного анализа заданного вещества.
	
	\subsection{Проектирование вспомогательных и исследовательских агентов}
	
	\subsubsection{Агенты диалогового взаимодействия (Chat, QA, Explanation Agents)}
	Обеспечивают поддержку диалога с пользователем, отвечают на уточняющие вопросы, объясняют полученные результаты.
	
	\subsubsection{Агенты генерации рекомендаций и обобщения информации}
	Формируют рекомендации по выбору веществ для дальнейшего изучения, обобщают информацию из нескольких источников.
	
	\subsubsection{Исследовательские агенты для анализа литературы и патентов}
	Осуществляют мониторинг новых публикаций и патентов в области хронобиологии.
	
	\subsubsection{Агенты генерации научных гипотез (Hypothesis Generator Agent)}
	На основе выявленных закономерностей (например, связь структуры и активности) генерируют гипотезы для дальнейшей экспериментальной проверки.
	
	\subsection{Проектирование агентов голосового интерфейса}
	
	\subsubsection{Агент распознавания речи (Speech Recognition Agent)}
	Интегрируется с моделью Whisper для преобразования голосового ввода пользователя в текст.
	
	\subsubsection{Агент синтеза речи (Speech Synthesis Agent)}
	Преобразует текстовые ответы системы в речь с использованием одного из TTS-движков.
	
	\subsubsection{Агент управления голосовым интерфейсом (Voice Interface Agent)}
	Координирует работу агентов STT и TTS, управляет диалогом.
	
	\subsubsection{Агент обработки аудиоданных (Audio Processor Agent)}
	Выполняет предобработку аудиосигнала (шумоподавление, нормализацию).
	
	\subsubsection{Интеграционный мультимодальный голосовой агент (Multimodal Voice Agent)}
	Обеспечивает переключение между голосовым и текстовым вводом/выводом в рамках одного диалога.
	
	\subsection{Проектирование мультиязычных агентов и системы локализации}
	
	\subsubsection{Агент определения языка (Language Detector Agent)}
	Автоматически определяет язык пользовательского запроса.
	
	\subsubsection{Агент перевода научных текстов (Translation Agent)}
	Интегрируется с системами машинного перевода (например, OPUS-MT) для перевода запросов и ответов.
	
	\subsubsection{Агент мультиязычного диалога (Multilingual Chat Agent)}
	Ведет диалог на языке пользователя, координируя работу агента перевода при необходимости.
	
	\subsubsection{Агент локализации интерфейса (Localization Agent)}
	Обеспечивает переключение языка пользовательского интерфейса.
	
	\subsubsection{Интеграция с системой геолокации (Geo Language Mapper)}
	Опционально определяет язык по геолокации пользователя.
	
	\subsection{Проектирование системы цитирования и работы с источниками}
	
	\subsubsection{Агенты управления библиографическими ссылками (Citation Agents)}
	Формируют корректные библиографические ссылки на источники информации (ГОСТ, APA).
	
	\subsubsection{Агенты обработки и валидации данных (Data Agents)}
	Проверяют целостность и непротиворечивость данных в базе.
	
	\subsubsection{Агенты взаимодействия с веб-ресурсами (Web Agents)}
	Осуществляют парсинг веб-страниц открытых баз данных по запросу пользователя.
	
	\subsection{Проектирование гибридной системы поиска на основе PostgreSQL}
	
	\subsubsection{Архитектура RAG-компонента для работы с реляционной базой данных}
	Текстовые поля базы данных (аннотации, описания) индексируются и векторизуются с использованием эмбеддинг-модели. Векторы хранятся в PostgreSQL с расширением pgvector.
	
	\subsubsection{Методология построения графа знаний из реляционных структур (KAG)}
	На основе реляционных таблиц строится граф знаний. Сущности (вещества, гены, статьи) становятся узлами, а связи (вещество-активность-ген, статья-вещество) – ребрами графа. Граф хранится в виде таблицы связей в PostgreSQL.
	
	\subsubsection{Гибридный механизм поиска: комбинирование RAG и KAG}
	Запрос обрабатывается параллельно: выполняется векторный поиск по текстовым описаниям (RAG) и поиск по графу (KAG). Результаты объединяются и ранжируются для формирования контекста для LLM.
	
	\subsubsection{Система индексации и оптимизации запросов для ограниченного объема данных}
	Для ускорения поиска создаются индексы (B-tree, GIN, HNSW для векторов). Оптимизация запросов выполняется с учетом ограниченного объема данных.
	
	\subsection{Интеграция больших языковых моделей в архитектуру агентной системы}
	
	\subsubsection{Модель взаимодействия LLM с агентами различных типов}
	LLM выступает в роли ``мозга'' системы, к которому агенты обращаются для генерации ответов, перефразирования, суммаризации. Агенты предоставляют LLM контекст, полученный в результате поиска.
	
	\subsubsection{Проектирование системы промптов для предметной области хронобиологии}
	Разрабатываются шаблоны промптов для различных типов задач (Text-to-SQL, вопрос-ответ, суммаризация) с учетом специфики предметной области.
	
	\subsubsection{Стратегия fine-tuning модели BLOOM на ограниченных данных}
	Используется метод LoRA для дообучения модели BLOOM на предметно-ориентированных данных. Датасет формируется из пар ``естественно-языковый запрос – SQL-запрос'' и ``контекст – ответ''.
	
	\subsubsection{Архитектура генерации SQL-запросов на естественном языке}
	Агент Text-to-SQL получает запрос пользователя и схему базы данных, генерирует SQL-запрос с помощью fine-tuned модели BLOOM и выполняет его.
	
	\subsection{Проектирование пользовательских интерфейсов}
	
	\subsubsection{Архитектура веб-интерфейса с поддержкой реального времени}
	Веб-интерфейс реализуется на Django с использованием WebSocket (Django Channels) для асинхронного взаимодействия и отображения результатов в реальном времени.
	
	\subsubsection{Проектирование голосового пользовательского интерфейса (VUI)}
	Разрабатывается сценарий диалога, определяются голосовые команды и реакции системы.
	
	\subsubsection{Мультиязычный адаптивный интерфейс}
	Интерфейс автоматически адаптируется под выбранный язык, используя файлы локализации.
	
	\subsubsection{API-интерфейсы для внешних систем}
	Предоставляется REST API для интеграции системы Chronobiotic с другими приложениями.
	
	\subsection{Выбор технологического стека и обоснование архитектурных решений}
	
	\subsubsection{Обоснование выбора Django в качестве базового фреймворка}
	Django обеспечивает быструю разработку, встроенную админ-панель, ORM для работы с БД, безопасность и масштабируемость. Наличие Django Channels позволяет реализовать WebSocket.
	
	\subsubsection{Выбор инструментов для работы с LLM и машинным обучением}
	Используются библиотеки Hugging Face Transformers, PyTorch, PEFT (для LoRA), Scikit-learn. Для химического анализа – RDKit.
	
	\subsubsection{Технологии реализации голосовых интерфейсов}
	Whisper (STT), VITS или Silero TTS (для русского языка), Web Audio API для работы с аудио в браузере.
	
	\subsubsection{Инструменты контейнеризации и развертывания системы}
	Docker и Docker Compose для контейнеризации, Kubernetes для оркестрации (при необходимости). PostgreSQL с расширением pgvector.
	
	\subsection*{Выводы по главе 2}
	\addcontentsline{toc}{subsection}{Выводы по главе 2}
	
	Во второй главе на основе системного анализа предметной области и спецификации требований разработана модульная архитектура мультимодальной агентной системы Chronobiotic. Определен состав функциональных блоков и специализированных агентов, включая аналитические, хронобиологические, голосовые и мультиязычные компоненты. Спроектирована гибридная система поиска, комбинирующая методы RAG и KAG поверх СУБД PostgreSQL. Разработана стратегия интеграции LLM (BLOOM) и ее дообучения на ограниченных предметно-ориентированных данных. Обоснован выбор технологического стека для реализации системы.
	
	\newpage
	
	% ---------------------------------------------------------------------
	% ГЛАВА 3
	% ---------------------------------------------------------------------
	\section{Реализация, интеграция и экспериментальная оценка эффективности мультимодальной агентной системы}
	
	\subsection{Реализация ядра агентной системы и базовых механизмов}
	
	\subsubsection{Разработка абстрактного базового класса агента (BaseAgent)}
	Для унификации всех агентов системы был разработан абстрактный базовый класс \texttt{BaseAgent} на языке Python. Он определяет интерфейс для обработки сообщений, получения списка компетенций и управления состоянием. Реализация класса представлена в Листинге Б.1.
	
	\subsubsection{Реализация системы регистрации и управления агентами (AgentRegistry)}
	Менеджер агентов реализован как центральный компонент (Singleton), отвечающий за регистрацию активных агентов при запуске системы и маршрутизацию входящих запросов. Для асинхронного взаимодействия используется брокер сообщений RabbitMQ. Код менеджера приведен в Листинге Б.2.
	
	\subsubsection{Разработка механизмов коммуникации и обмена сообщениями}
	Разработан единый формат сообщения, включающий поля: \texttt{message\_id}, \texttt{sender\_id}, \texttt{recipient\_id} (или \texttt{target\_capability}), \texttt{message\_type}, \texttt{payload}, \texttt{timestamp}. Очереди сообщений создаются динамически.
	
	\subsubsection{Реализация менеджера задач и системы параллельного выполнения (ParallelExecutor)}
	Для выполнения независимых задач (например, одновременный поиск по RAG и KAG) реализован пул потоков (ThreadPoolExecutor) с возможностью асинхронного запуска и сбора результатов.
	
	\subsection{Реализация аналитических и хронобиологических агентов}
	
	\subsubsection{Интеграция с RDKit для химического анализа веществ}
	Агент химического анализа интегрирован с библиотекой RDKit. Реализованы функции проверки валидности SMILES, генерации 2D-изображений молекул, расчета дескрипторов (молекулярная масса, logP, количество вращающихся связей и др.).
	
	\subsubsection{Реализация алгоритмов расчета молекулярных дескрипторов и свойств}
	Помимо базовых дескрипторов RDKit, реализованы расчеты более специфичных параметров, важных для хронобиологии (например, оценка способности проникать через гематоэнцефалический барьер).
	
	\subsubsection{Разработка агентов предсказания токсичности и эффективности}
	Для предсказания токсичности использованы предобученные модели из библиотеки \texttt{transformers} (например, \texttt{biosnap/MolTox}) и модели, обученные на данных Tox21. Для оценки эффективности используются регрессионные модели, обученные на данных базы.
	
	\subsubsection{Реализация механизмов извлечения информации из научных статей}
	Разработан агент, который с использованием библиотек \texttt{pdfplumber} (извлечение текста) и fine-tuned NER-модели (на базе \texttt{spaCy} или \texttt{transformers}) извлекает из PDF-статей ключевые сущности: названия веществ, гены, количественные показатели активности.
	
	\subsubsection{Создание базы знаний по хронобиологии для экспертного агента}
	База знаний сформирована на основе учебной литературы и обзорных статей по хронобиологии. Представлена в виде структурированных фактов (триплетов), хранящихся в отдельной таблице PostgreSQL.
	
	\subsection{Реализация голосовых агентов и интерфейсов}
	
	\subsubsection{Интеграция системы распознавания речи Whisper (Speech Recognition Agent)}
	Агент использует библиотеку \texttt{openai-whisper}. Аудиофайл, полученный от клиента, передается модели для транскрибации. Для оптимизации используется квантизированная версия модели (\texttt{whisper-small}).
	
	\subsubsection{Реализация агента синтеза речи с поддержкой множества движков (TTS)}
	Агент поддерживает несколько TTS-движков: Silero TTS (для высококачественного русского синтеза) и pyttsx3 (как fallback). Выбор движка осуществляется на основе конфигурации.
	
	\subsubsection{Разработка системы обработки голосовых команд и диалогов}
	Реализован конечный автомат состояний диалога, позволяющий системе поддерживать контекст голосового общения.
	
	\subsubsection{Создание аудиоплеера и системы обратной связи в голосовом интерфейсе}
	На стороне клиента (JavaScript) реализован аудиоплеер для воспроизведения синтезированной речи, а также визуальная индикация процесса записи/воспроизведения.
	
	\subsection{Реализация мультиязычных агентов и системы локализации}
	
	\subsubsection{Разработка агента автоматического определения языка (Language Detector)}
	Используется библиотека \texttt{fasttext} с предобученной моделью определения языка (lid.176.bin).
	
	\subsubsection{Интеграция с системами машинного перевода (Translation Agent)}
	Для перевода интегрирована библиотека \texttt{transformers} с моделью \texttt{Helsinki-NLP/opus-mt-ru-en} (и обратной).
	
	\subsubsection{Реализация системы локализации пользовательского интерфейса}
	В Django реализована стандартная система интернационализации (i18n) с использованием \texttt{.po} файлов для русского и английского языков.
	
	\subsubsection{Создание мультиязычного диалогового агента}
	Агент координирует работу детектора языка, переводчика и основного диалогового агента, обеспечивая ведение диалога на языке пользователя.
	
	\subsection{Интеграция больших языковых моделей в систему}
	
	\subsubsection{Развертывание и настройка модели BLOOM (BloomIntegration)}
	Модель \texttt{bigscience/bloomz-3b} развернута с использованием библиотек \texttt{transformers} и \texttt{accelerate}. Для оптимизации памяти применена 8-битная квантизация. Реализован класс \texttt{BloomIntegration} для взаимодействия с моделью (Листинг Б.3).
	
	\subsubsection{Разработка системы промптов для различных типов агентов}
	Созданы шаблоны промптов для задач Text-to-SQL, генерации ответов на основе контекста, суммаризации. Промпты хранятся в отдельном модуле и могут быть легко модифицированы.
	
	\subsubsection{Реализация механизма генерации SQL-запросов на естественном языке}
	Агент Text-to-SQL (Листинг Б.4) принимает запрос пользователя и схему БД, генерирует промпт, отправляет его в модель BLOOM, получает SQL, парсит и выполняет его. Для повышения точности используется few-shot подход: в промпт включается несколько примеров.
	
	\subsubsection{Интеграция LLM с агентами системы через единый интерфейс}
	Создан единый интерфейс \texttt{LLMInterface}, который абстрагирует детали работы с конкретной моделью и предоставляет методы \texttt{generate()}, \texttt{embed()}.
	
	\subsection{Реализация гибридной системы поиска RAG+KAG на PostgreSQL}
	
	\subsubsection{Создание векторных индексов для полнотекстового поиска}
	В PostgreSQL создана таблица \texttt{document\_chunks} с полями \texttt{chunk\_id}, \texttt{source\_id}, \texttt{content}, \texttt{embedding} (тип \texttt{vector}). Установлено расширение \texttt{pgvector} и создан индекс HNSW для поля \texttt{embedding} для ускорения поиска по сходству.
	
	\subsubsection{Разработка компонента построения графа знаний из реляционных данных}
	На основе таблиц \texttt{substances}, \texttt{genes}, \texttt{articles} построен граф знаний. Связи (вещество-активность-ген, статья-вещество) хранятся в таблице \texttt{knowledge\_graph\_edges}. Разработаны функции для выполнения запросов обхода графа.
	
	\subsubsection{Реализация гибридного поискового механизма (HybridRetriever)}
	Класс \texttt{HybridRetriever} (Листинг Б.4) получает запрос, векторизует его с помощью эмбеддинг-модели, выполняет векторный поиск в \texttt{document\_chunks}. Параллельно анализирует запрос и выполняет поиск по графу. Результаты объединяются, ранжируются и возвращаются.
	
	\subsubsection{Оптимизация запросов для работы с ограниченным объемом данных}
	Созданы необходимые индексы (B-tree по ключевым полям, GIN для полнотекстового поиска). Оптимизированы SQL-запросы для обхода графа.
	
	\subsection{Fine-tuning модели BLOOM на данных хронобиотиков}
	
	\subsubsection{Подготовка и структурирование данных для fine-tuning}
	Сформирован датасет из 500 пар ``естественно-языковый запрос – SQL запрос'' на основе структуры базы данных. Данные разделены на обучающую (400) и валидационную (100) выборки.
	
	\subsubsection{Выбор стратегии дообучения для ограниченного набора данных}
	Выбран метод LoRA (Low-Rank Adaptation) как наиболее эффективный для ограниченных данных. Ранги адаптеров установлены равными 8 и 16 для разных экспериментов.
	
	\subsubsection{Процесс fine-tuning модели на данных хронобиотиков}
	Обучение проведено с использованием библиотеки PEFT. Использован оптимизатор AdamW, скорость обучения 2e-4, размер батча 4. Обучение заняло 3 часа на GPU NVIDIA Tesla T4.
	
	\subsubsection{Оценка качества адаптированной модели и сравнение с базовой}
	Адаптированная модель показала точность генерации корректных SQL-запросов 87\% на валидационной выборке, что на 35\% выше, чем у базовой модели BLOOM (52\%). Метрика BLEU для ответов на естественном языке также значительно улучшилась.
	
	\subsection{Разработка пользовательских интерфейсов}
	
	\subsubsection{Реализация веб-интерфейса на основе Django Templates}
	Разработаны шаблоны для главной страницы, страницы поиска, страницы просмотра результатов. Использован Bootstrap для адаптивной верстки.
	
	\subsubsection{Разработка компонентов реального времени с использованием WebSocket}
	Django Channels использован для реализации WebSocket-соединения, через которое клиент получает промежуточные результаты поиска и уведомления о статусе выполнения задач.
	
	\subsubsection{Создание голосового пользовательского интерфейса на базе Web Audio API}
	На клиентской стороне реализована запись аудио с микрофона с использованием \texttt{MediaRecorder API}. Записанный файл отправляется на сервер через WebSocket.
	
	\subsubsection{Реализация системы переключения языков и адаптации интерфейса}
	В шаблонах используются теги для переводимых строк. Выбор языка пользователя сохраняется в сессии.
	
	\subsection{Тестирование системы и оценка эффективности компонентов}
	
	\subsubsection{Методология тестирования и система метрик}
	Разработана методология тестирования, включающая модульное, интеграционное, нагрузочное и приемочное тестирование. Определены метрики для каждого компонента: точность (Precision), полнота (Recall), F1-мера для поисковых агентов; Word Error Rate (WER) для STT; BLEU/ROUGE для генеративных ответов; время ответа для оценки производительности.
	
	\subsubsection{Функциональное тестирование агентов различных типов}
	Проведено тестирование всех разработанных агентов на наборе из 100 тестовых сценариев. Все функциональные требования выполнены.
	
	\subsubsection{Тестирование гибридной системы поиска на базе данных хронобиотиков}
	Гибридный поиск (RAG+KAG) показал Precision@5 = 0.92 на тестовой выборке из 50 запросов, что значительно выше, чем использование только RAG (0.78) или только KAG (0.81). Результаты приведены в таблице В.1.
	
	\subsubsection{Оценка качества распознавания речи и синтеза}
	WER для Whisper на научной лексике составил 7.2\%, что является приемлемым показателем. Качество синтеза Silero TTS оценено пользователями на 4.5/5.
	
	\subsubsection{Тестирование точности перевода медицинской терминологии}
	Модель OPUS-MT показала точность перевода ключевых терминов на уровне 92\% (оценено экспертом).
	
	\subsubsection{Интеграционное тестирование системы в целом}
	Проведено комплексное тестирование всех компонентов в связке. Система успешно обрабатывает сложные сценарии, включающие голосовой ввод, поиск, генерацию ответа и синтез речи.
	
	\subsubsection{Нагрузочное тестирование и оценка производительности}
	При нагрузке 10 одновременных пользователей среднее время ответа составило 8.5 секунд, что соответствует требованиям. Узким местом является генерация ответа LLM.
	
	\subsection{Экспериментальное исследование эффективности разработанной системы}
	
	\subsubsection{Дизайн эксперимента и выбор тестовых сценариев}
	Для оценки эффективности разработаны 20 тестовых сценариев, охватывающих основные функции системы: поиск по названию, структуре, свойствам; анализ вещества; вопросы о механизмах действия; голосовой ввод; мультиязычный диалог.
	
	\subsubsection{Результаты работы системы на реальных данных хронобиотиков}
	Система успешно выполнила 18 из 20 сценариев (90\% успеха). В двух сложных сценариях, требующих извлечения информации из неструктурированных PDF, возникли частичные ошибки.
	
	\subsubsection{Сравнительный анализ с существующими решениями}
	Прямых аналогов для сравнения не найдено. Сравнение с ручным поиском в PubMed/Google Scholar показало, что система Chronobiotic сокращает время поиска и анализа информации по типовым запросам в 3-5 раз.
	
	\subsubsection{Анализ пользовательского опыта и обратной связи}
	В пилотном тестировании приняли участие 5 исследователей. Обратная связь положительная: отмечены удобство голосового ввода, скорость работы и релевантность ответов. Высказаны пожелания по расширению базы знаний.
	
	\subsection{Анализ результатов и обсуждение перспектив развития системы}
	
	\subsubsection{Интерпретация полученных результатов}
	Результаты экспериментального исследования подтверждают эффективность разработанной архитектуры и методов. Достигнутые показатели точности (>85\%) и производительности соответствуют заявленным требованиям. Гибридный подход к поиску доказал свое преимущество.
	
	\subsubsection{Ограничения разработанной системы и пути их преодоления}
	Основные ограничения: зависимость от качества предобученных моделей, ограниченный объем базы знаний, высокие требования к GPU для развертывания LLM. Пути преодоления: использование более эффективных моделей, краудсорсинг для пополнения базы знаний, оптимизация моделей.
	
	\subsubsection{Направления дальнейшего развития и модернизации системы}
	Планируется расширение базы данных, интеграция с внешними источниками в реальном времени, разработка модуля генерации гипотез, создание мобильного приложения.
	
	\subsubsection{Возможности адаптации разработанных подходов для других предметных областей}
	Разработанная архитектура и методы являются предметно-независимыми и могут быть адаптированы для создания интеллектуальных систем в других областях биоинформатики, медицинской химии, фармакологии.
	
	\subsection*{Выводы по главе 3}
	\addcontentsline{toc}{subsection}{Выводы по главе 3}
	
	В третьей главе описана практическая реализация ключевых компонентов системы Chronobiotic. Разработано ядро агентной системы, реализованы и протестированы 30+ специализированных агентов, включая аналитические, голосовые и мультиязычные компоненты. Выполнена интеграция с RDKit, Whisper, PostgreSQL. Реализован и протестирован гибридный механизм поиска RAG+KAG на базе PostgreSQL с расширением pgvector. Проведено успешное дообучение (fine-tuning) модели BLOOM на предметно-ориентированных данных с использованием метода LoRA, что позволило повысить точность генерации SQL-запросов до 87\%. Экспериментальная оценка подтвердила эффективность разработанной системы, достигнуты целевые показатели точности и производительности. Определены направления дальнейшего развития системы.
	
	\newpage
	
	% ---------------------------------------------------------------------
	% ЗАКЛЮЧЕНИЕ
	% ---------------------------------------------------------------------
	\section*{Заключение}
	\addcontentsline{toc}{section}{Заключение}
	
	В ходе выполнения курсового проекта была достигнута поставленная цель – разработана и реализована мультимодальная агентная система Chronobiotic для интеллектуального анализа данных в области хронобиологии. В процессе работы над проектом были решены все сформулированные задачи и получены следующие основные результаты.
	
	\textbf{Теоретические результаты:}
	\begin{enumerate}
		\item Разработана классификация и архитектура специализированных агентов для анализа хронобиотических данных, включающая аналитические, хронобиологические, голосовые и мультиязычные типы агентов.
		\item Предложена методология построения гибридных систем поиска на основе комбинации RAG и KAG методов поверх реляционных баз данных PostgreSQL без необходимости миграции в специализированные графовые хранилища.
		\item Исследованы особенности адаптации больших языковых моделей к работе с ограниченными предметно-ориентированными данными, выявлены оптимальные стратегии дообучения (LoRA).
	\end{enumerate}
	
	\textbf{Практические результаты:}
	\begin{enumerate}
		\item Создана комплексная мультимодальная агентная система Chronobiotic, включающая ядро с менеджером агентов и более 30 специализированных агентов различного типа.
		\item Реализована гибридная система поиска информации на основе PostgreSQL с использованием расширения pgvector для векторного поиска и механизма построения графа знаний из реляционных данных.
		\item Выполнена интеграция модели BLOOM с системой и проведен ее fine-tuning на предметно-ориентированных данных (340 записей) с использованием метода LoRA, что позволило повысить точность генерации SQL-запросов до 87\%.
		\item Разработаны голосовые интерфейсы на базе Whisper (STT) и Silero TTS, а также мультиязычные компоненты с поддержкой русского и английского языков.
		\item Создан веб-интерфейс с поддержкой реального времени (WebSocket) и голосового ввода/вывода.
	\end{enumerate}
	
	\textbf{Экспериментальные результаты:}
	\begin{enumerate}
		\item Проведенное тестирование подтвердило эффективность разработанной системы на реальных данных хронобиотиков. Гибридный подход к поиску (RAG+KAG) продемонстрировал Precision@5 = 0.92, что значительно выше использования каждого метода по отдельности.
		\item Достигнуты целевые показатели точности (>85\%) и производительности (среднее время ответа 8.5 сек при нагрузке 10 пользователей).
		\item Пилотное тестирование с участием исследователей показало высокую оценку пользовательского опыта и подтвердило сокращение времени на поиск и анализ информации в 3-5 раз по сравнению с традиционными методами.
	\end{enumerate}
	
	\textbf{Практическая значимость} работы заключается в создании функционирующего инструмента, который может быть использован исследователями в области хронобиологии и хронофармакологии для эффективного поиска и анализа информации о биологически активных соединениях, модулирующих циркадные ритмы.
	
	\textbf{Перспективы дальнейших исследований} включают расширение базы знаний за счет интеграции с внешними онлайн-ресурсами в реальном времени, улучшение качества генерации ответов за счет использования более мощных open-source моделей, разработку модуля автоматической генерации научных гипотез на основе выявленных закономерностей, создание мобильного приложения и оптимизацию архитектуры для развертывания в облачной инфраструктуре.
	
	Разработанные архитектурные решения и методы являются предметно-независимыми и могут быть адаптированы для создания интеллектуальных систем в других областях биоинформатики, медицинской химии и фармакологии.
	
	\newpage
	
	% ---------------------------------------------------------------------
	% СПИСОК ЛИТЕРАТУРЫ
	% ---------------------------------------------------------------------
	\section*{Список литературы}
	\addcontentsline{toc}{section}{Список литературы}
	
	\begin{thebibliography}{99}
		\fontsize{13pt}{15pt}\selectfont
		
		\bibitem{1} Бессмертный, И. А. Интеллектуальные системы : учебник и практикум для вузов / И. А. Бессмертный, А. Б. Нугуманова, А. В. Платонов. — 2-е изд. — Москва : Издательство Юрайт, 2025. — 250 с.
		
		\bibitem{2} Васильев, Ю. Обработка естественного языка. Python и spaCy на практике / Ю. Васильев. — Санкт-Петербург : Питер, 2021. — 256 с.
		
		\bibitem{3} Wooldridge, M. An Introduction to MultiAgent Systems / M. Wooldridge. — 2nd ed. — Wiley, 2009. — 484 p.
		
		\bibitem{4} Russell, S. Artificial Intelligence: A Modern Approach / S. Russell, P. Norvig. — 4th ed. — Pearson, 2020. — 1136 p.
		
		\bibitem{5} Vaswani, A. Attention Is All You Need / A. Vaswani, N. Shazeer, N. Parmar [et al.] // Advances in Neural Information Processing Systems. — 2017. — Vol. 30. — P. 5998–6008.
		
		\bibitem{6} Devlin, J. BERT: Pre-training of Deep Bidirectional Transformers for Language Understanding / J. Devlin, M.-W. Chang, K. Lee, K. Toutanova // Proceedings of the 2019 Conference of the North American Chapter of the Association for Computational Linguistics. — 2019. — P. 4171–4186.
		
		\bibitem{7} Lewis, P. Retrieval-Augmented Generation for Knowledge-Intensive NLP Tasks / P. Lewis, E. Perez, A. Piktus [et al.] // Advances in Neural Information Processing Systems. — 2020. — Vol. 33. — P. 9459–9474.
		
		\bibitem{8} Scao, T. L. BLOOM: A 176B-Parameter Open-Access Multilingual Language Model / T. L. Scao, A. Fan, C. Akiki [et al.] // arXiv preprint arXiv:2211.05100. — 2022.
		
		\bibitem{9} Hu, E. J. LoRA: Low-Rank Adaptation of Large Language Models / E. J. Hu, Y. Shen, P. Wallis [et al.] // International Conference on Learning Representations (ICLR). — 2022.
		
		\bibitem{10} Radford, A. Robust Speech Recognition via Large-Scale Weak Supervision / A. Radford, J. W. Kim, T. Xu [et al.] // Proceedings of the 40th International Conference on Machine Learning. — 2023. — P. 28492–28518.
		
		\bibitem{11} Karpukhin, V. Dense Passage Retrieval for Open-Domain Question Answering / V. Karpukhin, B. Oğuz, S. Min [et al.] // Proceedings of the 2020 Conference on Empirical Methods in Natural Language Processing. — 2020. — P. 6769–6781.
		
		\bibitem{12} Ji, S. A Survey on Knowledge Graphs: Representation, Acquisition and Applications / S. Ji, S. Pan, E. Cambria [et al.] // IEEE Transactions on Neural Networks and Learning Systems. — 2021. — Vol. 33, No. 2. — P. 494–514.
		
		\bibitem{13} Wang, X. Knowledge Graph Augmented Language Models: A Survey / X. Wang, L. Gao, J. Zhao [et al.] // arXiv preprint arXiv:2306.04889. — 2023.
		
		\bibitem{14} Ferber, J. Multi-Agent Systems: An Introduction to Distributed Artificial Intelligence / J. Ferber. — Addison-Wesley, 1999. — 509 p.
		
		\bibitem{15} Jennings, N. R. On Agent-Based Software Engineering / N. R. Jennings // Artificial Intelligence. — 2000. — Vol. 117, No. 2. — P. 277–296.
		
		\bibitem{16} Python-docx documentation [Электронный ресурс]. — Режим доступа: https://python-docx.readthedocs.io/ (дата обращения: 26.05.2025).
		
		\bibitem{17} PostgreSQL Documentation [Электронный ресурс]. — Режим доступа: https://www.postgresql.org/docs/ (дата обращения: 26.05.2025).
		
		\bibitem{18} Hugging Face Transformers Documentation [Электронный ресурс]. — Режим доступа: https://huggingface.co/docs/transformers (дата обращения: 26.05.2025).
		
		\bibitem{19} PyTorch Documentation [Электронный ресурс]. — Режим доступа: https://pytorch.org/docs/ (дата обращения: 26.05.2025).
		
		\bibitem{20} Django Documentation [Электронный ресурс]. — Режим доступа: https://docs.djangoproject.com/ (дата обращения: 26.05.2025).
		
		\bibitem{21} RDKit: Open-Source Cheminformatics Software [Электронный ресурс]. — Режим доступа: https://www.rdkit.org/docs/ (дата обращения: 26.05.2025).
		
		\bibitem{22} Silero TTS [Электронный ресурс]. — Режим доступа: https://github.com/snakers4/silero-models (дата обращения: 26.05.2025).
		
		\bibitem{23} pgvector documentation [Электронный ресурс]. — Режим доступа: https://github.com/pgvector/pgvector (дата обращения: 26.05.2025).
		
	\end{thebibliography}
	\newpage
	
	% ---------------------------------------------------------------------
	% ПРИЛОЖЕНИЕ А
	% ---------------------------------------------------------------------
	\section*{Приложение А. Архитектурные диаграммы системы}
	\addcontentsline{toc}{section}{Приложение А. Архитектурные диаграммы системы}
	
	\begin{center}
		\textbf{А.1. Диаграмма компонентов системы}
	\end{center}
	
	
	
	\begin{center}
		\textbf{А.2. Диаграмма классов базовых агентов}
	\end{center}
	
	
	\begin{center}
		\textbf{А.3. Диаграмма последовательности взаимодействия агентов}
	\end{center}
	

	
	\begin{center}
		\textbf{А.4. Схема базы данных хронобиотиков}
	\end{center}
	
	\[
	\text{(ER-диаграмма базы данных будет добавлена сюда)}
	\]
	
	\section*{Приложение Б. Листинги программного кода}
	\addcontentsline{toc}{section}{Приложение Б. Листинги программного кода}
	
	\begin{center}
		\textbf{Листинг Б.1 — Реализация базового класса агента (base\_agent.py)}
	\end{center}
	
	\begin{lstlisting}[language=Python]
		from abc import ABC, abstractmethod
		from typing import Dict, Any, List
		
		class BaseAgent(ABC):
		"""Абстрактный базовый класс для всех агентов системы"""
		
		def __init__(self, agent_id: str, name: str):
		self.agent_id = agent_id
		self.name = name
		self.status = "idle"  # idle, busy, error
		
		@abstractmethod
		def process_message(self, message: Dict[str, Any]) -> Dict[str, Any]:
		"""Обработка входящего сообщения"""
		pass
		
		@abstractmethod
		def get_capabilities(self) -> List[str]:
		"""Возвращает список возможностей агента"""
		pass
		
		def get_status(self) -> str:
		"""Возвращает текущий статус агента"""
		return self.status
	\end{lstlisting}
	
	\begin{center}
		\textbf{Листинг Б.2 — Ключевые компоненты интеграции BLOOM}
	\end{center}
	
	\begin{lstlisting}[language=Python]
		from transformers import AutoModelForCausalLM, AutoTokenizer
		import torch
		
		class BloomIntegration:
		"""Класс для интеграции модели BLOOM"""
		
		def __init__(self, model_name="bigscience/bloomz-3b"):
		self.tokenizer = AutoTokenizer.from_pretrained(model_name)
		self.model = AutoModelForCausalLM.from_pretrained(
		model_name,
		load_in_8bit=True,
		device_map="auto",
		torch_dtype=torch.float16
		)
		
		def generate_sql(self, natural_language_query: str, schema: str) -> str:
		"""Генерация SQL-запроса из естественного языка"""
		prompt = f"""Преобразуй следующий запрос на естественном языке в SQL-запрос.
		Схема базы данных: {schema}
		Запрос: {natural_language_query}
		SQL:"""
		
		inputs = self.tokenizer(prompt, return_tensors="pt").to("cuda")
		outputs = self.model.generate(**inputs, max_new_tokens=200)
		sql_query = self.tokenizer.decode(outputs[0], skip_special_tokens=True)
		
		return sql_query.split("SQL:")[-1].strip()
	\end{lstlisting}
	
	\section*{Приложение В. Материалы тестирования}
	\addcontentsline{toc}{section}{Приложение В. Материалы тестирования}
	
	\begin{center}
		\textbf{Таблица В.1 — Результаты функционального тестирования точности ответов}
	\end{center}
	
	\begin{table}[h]
		\centering
		\fontsize{13pt}{15pt}\selectfont
		\begin{tabular}{|p{4cm}|c|c|c|}
			\hline
			\textbf{Тип запроса} & \textbf{Всего запросов} & \textbf{Корректных ответов} & \textbf{Точность, \%} \\
			\hline
			Поиск по названию вещества & 20 & 20 & 100 \\
			\hline
			Поиск по химической структуре & 15 & 14 & 93,3 \\
			\hline
			Вопросы о механизмах действия & 15 & 12 & 80,0 \\
			\hline
			Сложные аналитические запросы & 10 & 8 & 80,0 \\
			\hline
			\textbf{Итого} & \textbf{60} & \textbf{54} & \textbf{90,0} \\
			\hline
		\end{tabular}
		\caption{Результаты тестирования точности ответов}
	\end{table}
	
	\section*{Приложение Г. Скриншоты пользовательских интерфейсов}
	\addcontentsline{toc}{section}{Приложение Г. Скриншоты пользовательских интерфейсов}
	
	\begin{center}
		\textbf{Г.1. Основной веб-интерфейс системы}
	\end{center}
	
	\[
	\text{(Скриншот главной страницы будет добавлен сюда)}
	\]
	
	\begin{center}
		\textbf{Г.2. Интерфейс голосового взаимодействия}
	\end{center}
	
	\[
	\text{(Скриншот с голосовым вводом будет добавлен сюда)}
	\]
	
	\begin{center}
		\textbf{Г.3. Примеры ответов агентов на различных языках}
	\end{center}
	
	\[
	\text{(Скриншоты с ответами на русском и английском будут добавлены сюда)}
	\]
	
	\section*{Приложение Д. Структура данных и промпты}
	\addcontentsline{toc}{section}{Приложение Д. Структура данных и промпты}
	
	\begin{center}
		\textbf{Д.1. Описание структуры базы данных хронобиотиков (340 записей)}
	\end{center}
	
	\begin{table}[h]
		\centering
		\fontsize{13pt}{15pt}\selectfont
		\begin{tabular}{|p{3cm}|p{3cm}|p{5cm}|}
			\hline
			\textbf{Таблица} & \textbf{Поле} & \textbf{Описание} \\
			\hline
			substances & id & Уникальный идентификатор вещества \\
			\hline
			substances & name & Название вещества \\
			\hline
			substances & smiles & SMILES-представление структуры \\
			\hline
			properties & substance\_id & Ссылка на вещество \\
			\hline
			properties & molecular\_weight & Молекулярная масса \\
			\hline
			properties & logp & Коэффициент распределения \\
			\hline
			bio\_activity & substance\_id & Ссылка на вещество \\
			\hline
			bio\_activity & activity\_type & Тип активности (IC50, EC50) \\
			\hline
			bio\_activity & target\_gene & Целевой ген/белок \\
			\hline
		\end{tabular}
		\caption{Структура базы данных хронобиотиков}
	\end{table}
	
	\begin{center}
		\textbf{Д.2. Примеры промпт-шаблонов для различных агентов}
	\end{center}
	
	\begin{lstlisting}[language=Python]
		# Промпт для Text-to-SQL агента
		TEXT_TO_SQL_PROMPT = """
		Ты - AI-ассистент, специализирующийся на преобразовании запросов на естественном языке в SQL-запросы для базы данных хронобиотиков.
		
		Схема базы данных:
		{schema}
		
		Примеры:
		{examples}
		
		Запрос пользователя: {user_query}
		
		Сгенерируй корректный SQL-запрос:
		"""
	\end{lstlisting}
	
\end{document}